% =========================================================
% Proof: Half-Open Intervals Form a Ring
% Source: volume-i/algebras-of-sets/notes/...
% Mode: standard
% =========================================================

\subsection*{Half-Open Intervals Form a Ring}
\label{prf:half-open-intervals-ring}

\begin{remark*}[Return]
\hyperref[prop:half-open-intervals-ring]{$\leftarrow$ Back to Proposition (Half-Open Intervals Form a Ring) in Notes}
\end{remark*}

\begin{proof}
% TODO: proof to be filled.
% Strategy: Verify (R1)--(R3): difference of two half-open intervals is a finite union of half-open intervals.
\end{proof}

\begin{remark*}[Proof shape]
% TODO

\FlashProof{1. (R1): $\emptyset$ is the empty union, so $\emptyset\in\mathcal{H}$. 2. (R2): union of two finite unions of half-open intervals is again a finite union. 3. (R3): $(a,b]\setminus(c,d]$ splits into at most two half-open pieces at the endpoints; general case follows by finite case analysis.}
\end{remark*}

\begin{remark*}[Dependencies]
\begin{itemize}
  \item \hyperref[def:ring-of-sets]{Definition of ring of sets}: axioms (R1)--(R3)
\end{itemize}

\FlashUses{def:ring-of-sets}
\end{remark*}
